\documentclass[11pt]{article}
\usepackage[margin=1in]{geometry}
\usepackage{amsmath,amssymb,amsthm}
\usepackage[hidelinks]{hyperref}
\usepackage{xurl}
\setlength{\emergencystretch}{3em}
\newtheorem{theorem}{Theorem}
\newtheorem{lemma}[theorem]{Lemma}
\newtheorem{corollary}[theorem]{Corollary}
\newcommand{\Sig}{\Sigma_2}
\newcommand{\mr}{\operatorname{mr}}
\newcommand{\ar}{\operatorname{ar}}
\newcommand{\Vol}{\operatorname{Vol}}
\title{Nongraph double covers of a cofinite prime twist family}
\author{}
\date{17 November 2026}
\begin{document}
\maketitle
\begin{abstract}
For an explicit family of prime four-bridge knots with meridional rank
four and complex augmentation rank three, we prove that all but finitely
many double branched covers are not graph manifolds. The terminal twist
block gives fillings of one fixed torus-boundary manifold along pairwise
distinct slopes. A previously certified nonzero representation volume
for one filling excludes even connected sums of graph manifolds.
The reduced JSJ filling analysis of Boileau and Cooper then excludes
graph type for all but finitely many slopes. This gives
a qualitative cofinite conclusion, without an explicit cutoff or a
hyperbolicity assertion for the closed covers.
\end{abstract}

For $n\ge0$, let $K_n$ be the closure of the four-braid
\begin{equation}\label{eq:braid}
 \beta_n=\sigma_1\sigma_2\sigma_3^{-2}\sigma_2
 \sigma_3^{-2}\sigma_2\sigma_1^2\sigma_2^3
 \sigma_3^{-1}\sigma_2^{3+6n}.
\end{equation}
The independent family theorem \cite[Theorem \texttt{thm:main}]{family}
proves, for every $n\ge0$, that $K_n$ is prime and
\begin{equation}\label{eq:family}
 \mr(K_n)=b(K_n)=4,\qquad \ar(K_n,\mathbb C)=3,\qquad
 g(K_n)=7+3n,\qquad c(K_n)=17+6n.
\end{equation}
Here augmentation rank means ordinary complex $U=1$ augmentation
rank with meridian value different from one. This note concerns the
additional topology of their actual double branched covers.

We must distinguish irreducibility from graph structure. Call a compact
orientable manifold \emph{of graph type} when all its prime factors
are graph manifolds, allowing the $S^2\times S^1$ factors.
This is the connected-sum-closed graph convention used by Boileau and
Cooper \cite[printed p.~45]{bc}. For an irreducible manifold it is the
usual graph-manifold condition. Excluding graph type is stronger than
merely excluding irreducible graph manifolds.

\begin{theorem}\label{thm:cofinite}
There is a finite set $E\subset\mathbb Z_{\ge1}$ such that, for every
$n\ge0$ with $n\notin E$, the closed orientable manifold $\Sig(K_n)$
is not of graph type.
\end{theorem}

The point is that one nonzero representation-volume certificate can
control infinitely many different branch sets. Removing their common
terminal twist ball fixes a one-torus-boundary cover. If a hyperbolic
piece is separated from the boundary by graph pieces, good fillings
leave that piece intact. If the boundary reaches a hyperbolic piece
through a cable space, the cable filling either stays incompressible
or induces a Dehn filling on the hyperbolic piece. The exceptional
slopes in either case form a finite set.
The argument below is a consequence of the primary filling analysis
in \cite[Section~2]{bc}, with the graph-type and reducible cases made
explicit. It does not compute $E$, assert the conclusion for every
$n$, or establish hyperbolicity of the whole closed cover.

\section{One fixed tangle cover}

\begin{lemma}\label{lem:tangle}
There is a fixed compact, connected, orientable manifold $M$ with one
torus boundary and a basis $\lambda,\mu$ of $H_1(\partial M;\mathbb Z)$
such that
\[
 \Sig(K_n)\cong M(\alpha_n),\qquad
 \alpha_n=\lambda+(3+6n)\mu.
\]
In particular the unoriented primitive slopes $\alpha_n$ are pairwise
distinct, with
\[
                    \Delta(\alpha_n,\alpha_m)=6|n-m|.
\]
\end{lemma}
\begin{proof}
Choose a ball $B$ around the entire terminal two-strand twist block in
\eqref{eq:braid}, retaining its four marked boundary points.
Stretching the braid block inside $B$ accommodates every exponent
$3+6n$ without changing the complement tangle. The branch set in
$S^3\setminus\operatorname{int}B$ consists of two properly embedded
arcs: the knot has been cut at four points, and there are no remaining
closed components. Let $M$ be its double branched cover.
The character on the complement of the arcs sending a small meridian
of either arc to $1\in\mathbb Z/2$ is onto. The unbranched part of the
cover, and hence the branched cover, is connected. It is orientable.
Its boundary is the connected double cover of a sphere branched at
four points, so its Euler characteristic is $2\cdot2-4=0$ and it is
one torus.

The double cover of a rational two-string tangle is a solid torus.
For the trivial tangle this is the product of an annulus, the double
cover of a disk branched at two points, with an interval. Rational
tangles are obtained from it by marked-boundary homeomorphisms.
Gluing that solid torus to $M$ is Dehn filling. The rational-tangle
meridian formula and the distance formula are precisely those proved
in McCoy--Zentner \cite[Lemma~12 and its proof, p.~7]{mz}.
Their construction uses the cover of the fixed outside tangle, so the
same $M$ works simultaneously for the whole family, rather than a
potentially different parent for each pair of knots.

Choose $\lambda$ to be the disk slope of the zero-twist tangle.
The lift of one positive half-twist is a Dehn twist of the covering
annulus. Choose the sign of the primitive transverse class $\mu$
so that it sends $\lambda$ to $\lambda+\mu$; these two classes form
a basis. Equivalently, in the rational-tangle convention of
\cite[Figure~5 and Lemma~12]{mz}, integer twist tangles have rational
meridians affine in their integer parameter. Thus the block with $k$
positive half-twists has disk slope $\lambda+k\mu$. Applying this
with $k=3+6n$ proves the displayed slope formula.

The boundary gluing of the tangle pairs lifts to the boundary tori;
its two possible lifts differ by the covering involution and give the
same unoriented filled slope. The resulting double cover of $S^3$
is branched over the literal closure \eqref{eq:braid}. Uniqueness of
the double cover defined by the nontrivial meridional character
identifies it with $\Sig(K_n)$.
Finally the determinant of the coefficient vectors $(1,3+6n)$ and
$(1,3+6m)$ has absolute value $6|n-m|$. This proves distinctness
without choosing any software peripheral basis.
\end{proof}

\section{Nonzero volume and graph prime sums}

The stronger hypothesis needed from the initial knot is supplied by
its actual-cover certificate, not by a naming convention:
\begin{equation}\label{eq:seedvolume}
 \rho_0:\pi_1(\Sig(K_0))\longrightarrow\mathrm{PSL}_2(\mathbb C),
 \qquad 10.37<\Vol(\rho_0)<10.39.
\end{equation}
This is the closed representation constructed and certified in
\cite[Theorem \texttt{thm:nongraph} and its proof]{seed}; its filling
slope is tracked from the actual meridional double cover. We import
that nonzero-volume conclusion and do not replace it by a numerical
shape observation.

\begin{lemma}\label{lem:sumvolume}
Every $\mathrm{PSL}_2(\mathbb C)$ representation of a closed orientable
manifold of graph type has volume zero. Consequently $\Sig(K_0)$ is
not of graph type.
\end{lemma}
\begin{proof}
For a closed oriented manifold $Y$ and a representation $\rho$, its
volume is the integral over $Y$ of the pullback of the hyperbolic
volume form by any piecewise smooth equivariant map from its
universal cover to $\mathbb H^3$. The integral is independent of the
map \cite[Section~4.6, p.~11]{lty}.

If $Y=Y_1\#Y_2$, then $\pi_1Y=\pi_1Y_1*\pi_1Y_2$.
Let $\rho_i$ be the restrictions. Arrange an equivariant map to be
constant on each lift of the separating sphere collar. This is
possible because every lifted collar has trivial stabilizer, and
$\mathbb H^3$ is contractible; perform the modifications on orbit
representatives and translate equivariantly. On each punctured
summand cap its constant boundary sphere with a ball mapped
constantly. The integral over each collar and cap is zero. The
integral over the remaining pieces is exactly the sum of the
integrals defining the restricted volumes, so
\[
                 \Vol_Y(\rho)=\Vol_{Y_1}(\rho_1)+\Vol_{Y_2}(\rho_2).
\]
Iterate across a prime decomposition. On each irreducible graph
prime, every such volume vanishes by
\cite[Lemma~4.10, p.~11]{lty}. For $S^2\times S^1$, an equivariant
map may be chosen to factor through the lifted circle, so its
three-form integral is zero. Hence all terms vanish for graph type.
Equation~\eqref{eq:seedvolume} proves the final assertion.
\end{proof}

\section{Persistence under all but finitely many fillings}

We give the restricted consequence of Boileau--Cooper that is needed
here. In their terminology a \emph{hyperbolicabled} manifold is a
finite-volume hyperbolic core with disjoint cable spaces attached to
some boundary tori; each cable space has a free outer boundary torus.
Their reduced JSJ decomposition has blocks that are graph manifolds
or hyperbolicabled manifolds, with every cable subspace containing
an original boundary component of the whole manifold
\cite[pp.~50--51]{bc}.

\begin{lemma}\label{lem:persistence}
Let $P$ be a compact, connected, orientable, irreducible and
boundary-incompressible manifold with one torus boundary. If $P$ is
not a graph manifold, then $P(\alpha)$ is not of graph type for all
but finitely many slopes $\alpha$.
\end{lemma}
\begin{proof}
Take the reduced JSJ family of \cite[pp.~50--51]{bc}, and let $W$ be
the block incident to $\partial P$. Since $P$ is nongraph, some block
has a hyperbolic core. We explain why at least one such core survives
outside a finite set of slopes.

First suppose $W$ is a graph block. It cannot be a cable space:
the unique original boundary of $P$ belongs to $W$, so an adjacent
hyperbolicabled block has no cable subspace and is actually
hyperbolic. A single-cable $W$ could be absorbed into that neighbor.
A cable adjacent to a graph block could instead be combined with it.
Both operations contradict minimality. This is the argument in
\cite[proof of Lemma~2.3, p.~51]{bc}.
Nor can $W$ be $T^2\times I$. Its other boundary torus would either
be a second boundary component of $P$, impossible here, or would be
an interface which can be deleted by absorbing the product collar
into its neighbor, again contradicting minimality.
Apply Boileau--Cooper's Lemma~2.3 to $P$, not as an unrestricted
assertion about every graph block. Its hypotheses hold: $P$ is
compact, connected, orientable, irreducible and boundary-incompressible;
its boundary is one torus, and $P$ is not a cable space because it is
nongraph. The product $T^2\times I$ cannot be $P$ either.
The lemma gives a finite set of bad slopes on the sole boundary.
By its definition of a bad slope \cite[printed p.~51]{bc}, outside
this set the graph block $W(\alpha)$ has incompressible remaining
boundary tori. Only this incompressibility conclusion is needed here.
The preceding minimality argument explicitly excludes both the cable
and product forms for the reduced boundary block as well.

Outside this finite set, these interfaces remain incompressible after
filling and after attaching the unchanged blocks beyond $W$.
The subsequent argument does not require an irreducibility claim
for the filled graph block. Prime spheres can
be chosen disjoint from the incompressible interfaces by the usual
innermost-disk surgeries on their intersection circles. An unchanged
hyperbolic core beyond $W$ therefore survives in a prime factor,
separated by incompressible tori. Reducing the torus decomposition
can merge graph pieces or remove product collars, but cannot remove
this hyperbolic core. The filling is not of graph type.

Now suppose $W$ is hyperbolicabled, with hyperbolic core $H$.
There is at most one attached cable space, since every cable must
contain the unique original boundary component of $P$. All remaining
interfaces are therefore direct boundary tori of $H$.
We use the proof of \cite[Lemma~2.2, p.~50]{bc}, which retains the
hyperbolic core and gives more than just geometrizability.
If $\partial P$ is directly a boundary torus of $H$, the hyperbolic
Dehn-filling theorem used there excludes only finitely many slopes
and leaves a hyperbolic core, with all other cusps unfilled.
Otherwise $\partial P$ lies on a cable space $C$ attached to $H$.
Let $f$ be its regular fiber slope and let $s\ge2$ be its exceptional
fiber order. Explicitly exclude the single slope $f$, for which
$\Delta(\alpha,f)=0$.
If $\Delta(\alpha,f)\ge2$, then $C(\alpha)$ is Seifert fibered over a
disk with two exceptional fibers and incompressible boundary. Thus
$H$ remains separated by incompressible tori and is unchanged.
If $\Delta(\alpha,f)=1$, then $C(\alpha)$ is a solid torus inducing a
slope $\gamma(\alpha)$ on the adjacent boundary of $H$. The cable
calculation in \cite[proof of Lemma~2.2, p.~50]{bc} gives
\[
 \Delta\bigl(\gamma(\alpha),\gamma(\alpha')\bigr)
       =s^2\Delta(\alpha,\alpha').
\]
Hence this map of slopes is injective. Only finitely many $\alpha$
can induce one of the finitely many exceptional hyperbolic filling
slopes of $H$. For every remaining slope, $H$ either stays unchanged
or undergoes a hyperbolic filling. The incompressible remaining
interfaces preserve a hyperbolic JSJ piece in $P(\alpha)$.
This proves the lemma.
\end{proof}

The wording of Lemma~\ref{lem:persistence} is deliberately restricted.
The conclusion is read from the hyperbolic-core argument in the
proofs of Boileau--Cooper's Lemmas~2.2--2.3, not from the word
``geometrizable'' alone. Also the product block $T^2\times I$ was
excluded explicitly. The indispensable finite-bad-slope import is the restricted
one-boundary nongraph application of Boileau--Cooper just made.
The original boundary-compression theorem of Culler--Gordon--Luecke--Shalen
\cite[Theorem~2.4.4, printed p.~276]{cgls} explicitly displays cable
and product exceptions and explains the need for this distinction.
That original statement is context, not a separate premise of the
proof here. No all-good-filling irreducibility claim is needed in the
graph-block case.

\begin{theorem}\label{thm:onefilling}
Let $M$ be any compact, connected, orientable manifold with one torus
boundary. If one Dehn filling of $M$ is not of graph type, then all
but finitely many slopes give fillings not of graph type.
\end{theorem}
\begin{proof}
Use a prime decomposition with the boundary in the factor $P$:
\[
                     M=P\#Q_1\#\cdots\#Q_r,
\]
where the $Q_i$ are closed prime factors. For every slope $\alpha$,
\[
                 M(\alpha)=P(\alpha)\#Q_1\#\cdots\#Q_r.
\]
If any $Q_i$ is not a graph prime, it survives every filling, so the
conclusion holds for every slope. Assume otherwise that all $Q_i$
are graph primes. A graph manifold remains of graph type after
solid-torus filling. Indeed, one can append the solid torus to its
Seifert-piece gluing description; the gluing tori need not be
incompressible. Equivalently this is graph-filling closure in the
connected-sum convention of \cite[printed pp.~45 and~51]{bc}.
Thus if $P$ were graph, every filling of $M$ would be of graph type,
contrary to the hypothesis.

The prime factor $P$ has nonempty torus boundary, so it is
irreducible: the exceptional orientable prime $S^2\times S^1$ is
closed. It is also boundary-incompressible. An irreducible manifold
with compressible torus boundary is a solid torus, which would be
graph. Therefore Lemma~\ref{lem:persistence} applies to $P$.
For every slope outside its finite exceptional set, $P(\alpha)$ has
a hyperbolic JSJ piece. Adding the fixed graph prime factors does
not remove that piece, proving the theorem.
\end{proof}

\section{Application to the prime knots}

\begin{proof}[Proof of Theorem~\ref{thm:cofinite}]
By Lemma~\ref{lem:tangle}, all actual branched covers are fillings of
one $M$ along the distinct slopes $\alpha_n$.
Lemma~\ref{lem:sumvolume} and the imported nonzero volume
\eqref{eq:seedvolume} show that $M(\alpha_0)$ is not of graph type.
Theorem~\ref{thm:onefilling} leaves only finitely many slopes which
could give graph type. Their intersection with the injective sequence
$\{\alpha_n:n\ge0\}$ has finitely many indices. Removing $0$, whose
cover is already excluded, yields the required finite set $E$.
\end{proof}

Combining this with \eqref{eq:family} gives infinitely many pairwise
distinct prime knots with $\ar(K_n,\mathbb C)=3<4=\mr(K_n)=b(K_n)$
whose actual double branched covers are outside graph type. This
excludes any equality-family hypothesis requiring a graph double
cover for these cofinite members. It is not an exclusion of every
previously known meridional-rank equality theorem.

The same cover exclusion persists after adding trefoils. The double
branched cover of a connected sum is the connected sum of the covers:
the sum sphere meets the branch knot twice, so its lift is a sphere.
Thus, for $n\notin E$ and every $s\ge0$,
\[
 \Sig\bigl(K_n\#(\#^s3_1)\bigr)
       \cong\Sig(K_n)\#\bigl(\#^s\Sig(3_1)\bigr)
\]
still has a nongraph prime factor. The rank conclusions already
proved in \cite[Corollary \texttt{cor:sums}]{family} remain
$\ar=3+s$ and $\mr=b=4+s$; no rank additivity is inferred merely
from this topological connected-sum observation.

\paragraph{Dependencies and scope.}
The new argument is qualitative and topological. It imports the
literal twist-family definition and its independently proved rank,
primeness and size results, the initial actual-cover nonzero-volume
certificate, the rational-tangle filling construction, graph-prime
representation-volume vanishing, and the restricted reduced-JSJ
filling analysis. No new triangulation or sampled parameter value is
needed as a premise. The source audit retains exact editions and
locators; it is a source and proof record, not a claimed numerical
or formal verification. There is no explicit bound on $E$, no claim
that $E$ is empty, and no assertion that all these closed covers are
hyperbolic.

\paragraph{AI assistance and review.}
ChatGPT (OpenAI) assisted with mathematical arguments, literature checks,
exposition, proof review and verification code.

\begin{thebibliography}{9}\small\raggedright
\bibitem{bc} M.~Boileau and D.~Cooper,
\emph{On three-manifolds with bounded geometry},
Bol.~Soc.~Mat.~Mexicana (3) 10, Special Issue (2004), 43--53.
Published author-hosted copy,
\url{https://web.math.ucsb.edu/~cooper/40.pdf}.
Graph convention p.~45; Lemma~2.2 and proof p.~50;
reduced JSJ family, Lemma~2.3 and proof p.~51.
\bibitem{cgls} M.~Culler, C.~McA.~Gordon, J.~Luecke and P.~B.~Shalen,
\emph{Dehn surgery on knots}, Ann.~of Math.~(2) 125 (1987), 237--300.
Section~2.4 setup, printed pp.~274--275, and Theorem~2.4.4,
printed p.~276 (PDF p.~41),
\url{https://marc-culler.info/static/home/papers/CyclicSurgery.pdf}.
\bibitem{mz} D.~McCoy and R.~Zentner,
\emph{The Montesinos trick for proper rational tangle replacement},
Proc.~Amer.~Math.~Soc.~151 (2023), no.~4, 1811--1822,
\href{https://doi.org/10.1090/proc/16246}{doi:10.1090/proc/16246}.
Retained author copy, Figure~5 p.~6, Lemma~12 and proof p.~7,
\url{https://zentner.app.uni-regensburg.de/proper_tangle_replace.pdf}.
\bibitem{lty} F.~Luo, S.~Tillmann and T.~Yang,
\emph{Thurston's spinning construction and solutions to the hyperbolic
gluing equations for closed hyperbolic 3-manifolds},
retained arXiv:1004.2992v1. Section~4.6 and Lemma~4.10, p.~11.
\bibitem{seed} \emph{An explicit seventeen-crossing four-bridge knot},
companion \path{NEW_SEVENTEEN_CROSSING_KNOT.tex},
Theorem \texttt{thm:nongraph} and its proof: actual closed-cover
representation with volume in $(10.37,10.39)$.
\bibitem{family} \emph{Prime four-bridge knots with augmentation rank
three}, companion \path{PRIME_AUGMENTATION_GAP_FAMILY.tex},
Theorem \texttt{thm:main}, Corollary \texttt{cor:sums}.
\end{thebibliography}
\end{document}
